% Part: computability
% Chapter: computability-theory
% Section: non-comp-set

\documentclass[../../../include/open-logic-section]{subfiles}

\begin{document}

\olfileid{cmp}{thy}{ncp}
\olsection{توجد مجموعات غير قابلة للحساب}

قد ثبت أن قابلية المجموعة للحساب تستلزم قابليتها للتعداد حسابيًا. وأما
العكس فلا يصح على العموم، كما تبيّن النتيجة الآتية.

\begin{thm}\ollabel{thm:K-0}
لتكن $K_0$ المجموعة $\Setabs{\tuple{e, x}}{\cfind{e}(x) \fdefined}$.
عندئذ تكون $K_0$ قابلة للتعداد حسابيًا، ولكنها غير قابلة للحساب.
\end{thm}

\begin{proof}
أما قابلية $K_0$ للتعداد حسابيًا فتثبت بجعلها مجالًا للدالة~$f$
التي تعريفها
\[
f(z) = \umin{y}{(\len{z} = 2 \land T((z)_0, (z)_1, y))}.
\]
إذ متى كانت $\cfind{e}(x)$ معرّفة عثرت $f(\tuple{e,x})$ على سجل حساب
متوقف؛ وإن لم تُعرّف $\cfind{e}(x)$، لم تُعرّف
$f(\tuple{e,x})$. وأما إن لم يرمز $z$ إلى زوج، فـ$f(z)$
غير معرّفة كذلك.

وأما انتفاء قابلية $K_0$ للحساب فهو عين عدم قابلية مشكلة التوقف
للقرار في \olref[hlt]{thm:halting-problem}.
\end{proof}

فـ$K_0$ تجمع الأزواج $\tuple{e,x}$ التي تحقق
$\cfind{e}(x)\fdefined$؛ ومعنى $\tuple{e,x}\in K_0$ أن
$\cfind{e}$ تتعين عند الدخل~$x$، أي يتوقف حسابها هناك. ومن ثم تسمى
«مجموعة التوقف». ونخص المجموعة
$K=\Setabs{e}{\cfind{e}(e)\fdefined}$ باسم «مجموعة التوقف الذاتي»؛
وهي مثال معياري كثير الاستعمال لما لا يقبل القرار.

\begin{thm}\ollabel{thm:K}
مجموعة التوقف الذاتي $K=\Setabs{e}{\cfind{e}(e)\fdefined}$
!!{c.e.} ولكنها غير قابلة للقرار.
\end{thm}

\begin{proof}
  لنفرض، طلبًا للخلف، قابلية $K$ للقرار؛ فتكون دالتها المميزة $\Char{K}$ قابلة
  للحساب. ونضع
  \[d(e) = \begin{cases}
  1 & \text{إذا كانت\/ $\Char{K}(e) = 0$}\\
  \fundefined & \text{في غير ذلك.}
  \end{cases}
  \]
  ثم نأخذ $k$ فهرسًا لـ~$d$، فيكون $d\simeq\cfind{k}$، ومنه
  $d(k)\simeq\cfind{k}(k)$. وهذا يناقض حقيقة أن
  $d(k)\fdefined$ إذا وفقط إذا كانت $\cfind{k}(k)\fundefined$، كما
  يقتضي تعريف~$d$.

  وأما $K$ فهي مجال $f(x)=\umin{y}{T(x,x,y)}$، ومن ثم فهي !!{c.e.}.
\end{proof}

\end{document}
